\chapter{连续性}
\label{chap:continuity-discontinuity}

当函数在一点的取值等于它在该点的极限时，局部极限便与函数本身接合起来。连续性首先是逐点性质；加入单调性以后，左右极限总会存在，间断只能以跳跃方式出现，逆映射也获得相应的连续性。

\section{点连续、序列刻画与相对连续}

\subsection{点连续的定义及极限形式}


\begin{definition}[点连续]
设 \(f:D\to\R\)，\(a\in D\)。若
\[
 \forall\varepsilon>0\ \exists\delta>0\ \forall x\in D,\qquad
 \abs{x-a}<\delta\Longrightarrow\abs{f(x)-f(a)}<\varepsilon,
\]
则称 \(f\) 在 \(a\) 处连续。
\end{definition}

与删心极限相比，连续性允许 \(x=a\)，但此时误差自动为零。若 \(a\) 是 \(D\) 的聚点，则
\[
 f\text{ 在 }a\text{ 连续}
 \quad\Longleftrightarrow\quad
 \lim_{\substack{x\to a\\x\in D}}f(x)=f(a).
\]
若 \(a\) 是孤立点，任意函数都在 \(a\) 连续，因为可取小邻域使其中唯一的定义域点是 \(a\)。

点连续性比较的是两类同时变化的数据：输入点 \(x\) 靠近 \(a\)，输出值 \(f(x)\) 靠近已经指定的 \(f(a)\)。因此连续性不仅要求删心极限存在，还要求点值恰好填在这个极限上。可去间断点之所以能修补，正是因为邻近函数值已经决定了唯一可选的点值。

\begin{definition}[集合上的连续性]
若 \(f\) 在 \(D\) 的每一点连续，则称 \(f\) 在 \(D\) 上连续。若 \(D\) 是区间，端点连续性始终相对于 \(D\) 理解。
\end{definition}

\begin{example}[可去补值]
设 \(f(x)=(x^2-1)/(x-1)\)（\(x\ne1\)）。因删心极限为 \(2\)，定义 \(f(1)=2\) 可得到连续延拓；定义为其他值则在 \(1\) 处不连续。
\end{example}

\subsection{连续性的序列刻画}


\begin{theorem}[序列连续性]
设 \(a\in D\)。函数 \(f:D\to\R\) 在 \(a\) 连续，当且仅当对每条 \((x_n)\subset D\) 满足 \(x_n\to a\)，都有
\[
 f(x_n)\to f(a).
\]
\end{theorem}

\begin{proof}
若 \(f\) 连续，给定 \(\varepsilon>0\)，先取连续性半径 \(\delta\)，再用 \(x_n\to a\) 使 \(\abs{x_n-a}<\delta\)，即得像列收敛。

反之，若 \(f\) 在 \(a\) 不连续，则存在 \(\varepsilon_0>0\)，使对每个 \(n\) 可选 \(x_n\in D\) 满足
\[
 \abs{x_n-a}<1/n,\qquad
 \abs{f(x_n)-f(a)}\ge\varepsilon_0.
\]
于是 \(x_n\to a\)，但 \(f(x_n)\not\to f(a)\)，与序列条件矛盾。
\end{proof}

\begin{corollary}
连续函数把收敛数列映为收敛数列，并保持极限：
\[
 x_n\to a\in D\Longrightarrow f(x_n)\to f(a).
\]
\end{corollary}

条件 \(a\in D\) 不可省略，因为 \(f(a)\) 必须有定义。若只研究 \(a\notin D\) 的邻近行为，应使用函数极限而非点连续。

\subsection{相对拓扑下的端点连续性}

\begin{definition}[开集、闭集与相对开集]
集合 \(U\subset\R\) 称为\term{开集}，如果对每个 \(x\in U\)，都存在 \(r>0\) 使
\[
 (x-r,x+r)\subset U.
\]
集合 \(F\subset\R\) 称为\term{闭集}，如果 \(\R\setminus F\) 为开集。给定 \(D\subset\R\)，若 \(A=D\cap U\)，其中 \(U\) 为实线中的开集，则称 \(A\) 在 \(D\) 中\term{相对开}；若 \(D\setminus A\) 在 \(D\) 中相对开，则称 \(A\) 在 \(D\) 中\term{相对闭}。
\end{definition}

开区间、空集和 \(\R\) 都是开集；闭区间和单点集是闭集。半闭区间 \([a,b)\) 在 \(\R\) 中既不开也不闭，但它在 \([a,b]\) 中相对开，因为
\[
 [a,b)=[a,b]\cap(-\infty,b).
\]
“相对”二字记录了讨论邻近点时所采用的定义域。

\begin{theorem}[连续性的开集刻画]
函数 \(f:D\to\R\) 在 \(D\) 上连续，当且仅当对每个开集 \(V\subset\R\)，原像 \(f^{-1}(V)\) 在 \(D\) 中相对开。等价地，实线中每个闭集的原像在 \(D\) 中相对闭。
\end{theorem}

\begin{proof}
设 \(f\) 连续，取 \(x\in f^{-1}(V)\)。由 \(V\) 的开性，存在 \(\varepsilon>0\) 使
\((f(x)-\varepsilon,f(x)+\varepsilon)\subset V\)。连续性给出 \(\delta>0\)，使
\[
 y\in D,\quad \abs{y-x}<\delta
 \Longrightarrow f(y)\in V.
\]
故 \(f^{-1}(V)\) 在 \(D\) 中相对开。

反之，给定 \(x\in D\) 与 \(\varepsilon>0\)，区间
\(V=(f(x)-\varepsilon,f(x)+\varepsilon)\) 的原像在 \(D\) 中相对开且含 \(x\)，所以存在 \(\delta>0\) 使
\[
 (x-\delta,x+\delta)\cap D\subset f^{-1}(V).
\]
这正是 \(f\) 在 \(x\) 处连续。闭集版本对原像和补集使用
\(f^{-1}(\R\setminus F)=D\setminus f^{-1}(F)\) 即得。
\end{proof}


若 \(f:[a,b]\to\R\)，则在左端点 \(a\) 连续的量词条件为
\[
 \forall\varepsilon>0\ \exists\delta>0,\qquad
 x\in[a,b],\ a\le x<a+\delta
 \Longrightarrow\abs{f(x)-f(a)}<\varepsilon.
\]
它等价于右极限 \(\lim_{x\to a+}f(x)=f(a)\)。右端点相应使用左极限。

\begin{example}
\(f(x)=\sqrt x\) 在定义域 \([0,\infty)\) 的端点 \(0\) 连续，因为给定 \(\varepsilon>0\)，取 \(\delta=\varepsilon^2\)，则
\[
 0\le x<\delta\Longrightarrow\abs{\sqrt x-\sqrt0}<\varepsilon.
\]
没有必要把函数延拓到负半轴后再讨论。
\end{example}

\begin{proposition}[限制与粘合]
连续函数限制到任意子集仍连续。设 \(D=A\cup B\)，连续函数
\[
 f_A:A\to\R,\qquad f_B:B\to\R
\]
在 \(A\cap B\) 上取值相同。如果 \(A,B\) 在 \(D\) 中都相对开，或者都相对闭，则存在唯一连续函数 \(f:D\to\R\)，其在 \(A,B\) 上的限制分别为 \(f_A,f_B\)。
\end{proposition}

\begin{proof}
限制的连续性由定义直接得到。两个公式在交集上一致，故确实定义唯一函数 \(f\)。若 \(A,B\) 相对开，任取开集 \(V\subset\R\)，有
\[
 f^{-1}(V)=f_A^{-1}(V)\cup f_B^{-1}(V).
\]
右边两项分别在 \(A,B\) 中相对开；因 \(A,B\) 本身在 \(D\) 中相对开，它们也在 \(D\) 中相对开。因此 \(f^{-1}(V)\) 相对开，\(f\) 连续。

若 \(A,B\) 相对闭，则对任意闭集 \(F\subset\R\) 使用同一等式。右边两项分别是 \(A,B\) 中的相对闭集，因而是 \(D\) 中的相对闭集；有限并仍相对闭。闭集原像刻画再次给出连续性。
\end{proof}

“分段公式各自连续”不足以自动推出分段函数连续；必须另行核对交界点的函数值与两侧极限。

\section{连续映射的代数运算和复合}

连续性的运算法则不是另一套需要重新记忆的极限公式。只要把每个函数在 \(a\) 处的极限替换为它的点值，函数极限的加、乘、商和复合定理便直接给出连续性。商仍须排除分母点值为零；复合仍须保证内层函数值落在外层函数的定义域中。


\begin{theorem}
若 \(f,g:D\to\R\) 在 \(a\in D\) 连续，\(c\in\R\)，则
\[
 f+g,\quad cf,\quad fg,\quad\abs f,\quad
 \max\{f,g\},\quad\min\{f,g\}
\]
均在 \(a\) 连续。若 \(g(a)\ne0\)，则 \(f/g\) 在 \(a\) 连续。
\end{theorem}

\begin{proof}
把函数极限的代数法则与
\[
 \lim_{x\to a}f(x)=f(a),\qquad
 \lim_{x\to a}g(x)=g(a)
\]
结合即可。最大值与最小值使用
\[
 \max\{u,v\}=\frac{u+v+\abs{u-v}}2,\qquad
 \min\{u,v\}=\frac{u+v-\abs{u-v}}2.
\]
商的非零条件保证 \(g\) 在 \(a\) 附近非零。
\end{proof}

\begin{theorem}[复合连续性]
若 \(g:D\to E\) 在 \(a\) 连续，\(f:E\to\R\) 在 \(g(a)\) 连续，则 \(f\circ g\) 在 \(a\) 连续。
\end{theorem}

\begin{proof}
若 \(x_n\to a\)，则 \(g(x_n)\to g(a)\)，继而
\[
 f(g(x_n))\to f(g(a)).
\]
由序列连续性，\(f\circ g\) 在 \(a\) 连续。
\end{proof}

\begin{corollary}
多项式在 \(\R\) 上连续；有理函数在分母非零处连续；根式、绝对值及由有限次代数运算和复合得到的初等表达式在其自然定义域内连续。
\end{corollary}

\section{间断点的分类}


设 \(a\) 属于函数的定义域，并假定两侧都是定义域的聚集方向。
\begin{definition}[间断点分类]
\begin{enumerate}
 \item 若左右有限极限存在且相等为 \(L\)，但 \(f(a)\ne L\)，则称 \(a\) 为\term{可去间断点}；
 \item 若左右有限极限都存在但不相等，则称 \(a\) 为\term{跳跃间断点}；
 \item 若至少一个单侧有限极限不存在，则称 \(a\) 为\term{第二类间断点}。
\end{enumerate}
\end{definition}

若 \(a\) 不属于定义域，但函数在 \(a\) 的删心邻域内有有限极限 \(L\)，则应说函数可在缺点 \(a\) 处连续延拓：补定义 \(f(a)=L\) 后得到连续函数。它与定义域内因点值错误造成的可去间断使用同一个修补方法，但在术语上不是原函数定义域内的间断点。

分类只描述单侧极限怎样失败，并不衡量间断“大小”。可去间断可以只由一个错误点值造成，跳跃间断保留两个有限侧极限，而第二类间断把无界发散和持续振荡都包括在内。若定义域在 \(a\) 的一侧没有点，则那一侧极限不应被强行加入分类。

\begin{example}
把
\[
 f(x)=\frac{x^2-1}{x-1}\quad(x\ne1),\qquad f(1)=0
\]
看作定义在 \(\R\) 上的函数，则 \(1\) 是可去间断点；若不补上 \(x=1\)，原表达式只是在该缺点处可连续延拓。取整函数在每个整数处具有跳跃间断；\(\sin(1/x)\) 在补定义 \(f(0)=0\) 后，\(0\) 是第二类间断点。函数 \(1/x\) 在 \(0\) 处的左右扩充实极限存在但不是有限数，若在 \(0\) 任意补定义，按上述有限极限分类也归入第二类。
\end{example}

分类依赖定义域的趋近方向。在半闭区间端点只检验存在的单侧方向，因此不能机械套用双侧分类。

\section{单调函数的左右极限与间断点}

\subsection{单调函数的左右极限}


\begin{theorem}[单调函数的单侧极限]
设 \(f:I\to\R\) 在区间 \(I\) 上单调递增，\(a\) 是 \(I\) 的内点。则有限左右极限存在，并且
\[
 f(a-)=\sup\{f(x):x\in I,\ x<a\},
\]
\[
 f(a+)=\inf\{f(x):x\in I,\ x>a\}.
\]
此外
\[
 f(a-)\le f(a)\le f(a+).
\]
\end{theorem}

\begin{proof}
记左侧值集上确界为 \(L_-\)。它非空，且由任取 \(y>a\) 有 \(f(x)\le f(y)\)（\(x<a\)），故有上界。给定 \(\varepsilon>0\)，按上确界定义取 \(x_0<a\) 使
\[
 L_--\varepsilon<f(x_0)\le L_-.
\]
当 \(x_0<x<a\) 时，单调性给出
\[
 L_--\varepsilon<f(x_0)\le f(x)\le L_-,
\]
故 \(f(x)\to L_-\)（\(x\to a-\)）。右侧用下确界同理。最后的不等式直接来自单调性并取上确界、下确界。
\end{proof}

递减函数有对称结论，上确界与下确界的位置互换。区间端点若存在相应侧聚集，也可得到单侧极限，有限性由区间内任一另一点的函数值控制。

\begin{corollary}
单调函数的每个内点间断都是跳跃间断；它在 \(a\) 连续当且仅当
\[
 f(a-)=f(a)=f(a+).
\]
\end{corollary}

\subsection{单调函数的间断点至多可数}


\begin{theorem}
区间上的单调函数的间断点集合至多可数。
\end{theorem}

\begin{proof}
只证递增情形。对每个内点间断 \(a\)，上一节给出
\[
 f(a-)<f(a+).
\]
从非空开区间 \(J_a=(f(a-),f(a+))\) 中选一个有理数 \(q_a\)。

若 \(a<b\)，则对任意 \(x,y\) 满足 \(a<x<y<b\)，有 \(f(x)\le f(y)\)。令 \(x\to a+\)、\(y\to b-\)，得到
\[
 f(a+)\le f(b-).
\]
因此 \(J_a\) 与 \(J_b\) 不相交，所选有理数不同。映射 \(a\mapsto q_a\) 是从间断点集到 \(\Q\) 的单射，故内点间断至多可数。区间端点至多增加两个点。
\end{proof}

\begin{corollary}
单调函数在区间上除至多可数个点外处处连续。
\end{corollary}

“至多可数”不能加强为“有限”。例如定义 \(f:[0,1]\to\R\) 为
\[
 f(0)=0,
 \qquad
 f(x)=\frac1n\quad\left(\frac1{n+1}\le x<\frac1n,\ n\ge1\right),
 \qquad
 f(1)=1.
\]
当 \(x\) 增大时，相应的 \(n\) 只会减小，故 \(f\) 单调递增；它在每个 \(1/n\)（\(n\ge2\)）处都有跳跃。因此单调函数可以有无穷多个间断点。

\section{严格单调映射及其相对连续逆}

严格单调保证每个函数值只有一个原像，但并不保证像集没有缺口。逆函数的连续性应先相对于像集 \(J=f(I)\) 来讨论；在这个相对意义下，严格单调已经足够。要进一步断言 \(J\) 本身是区间，才需要连续性和下一章的介值定理。


\begin{theorem}[连续逆映射]
设 \(I\) 是区间，\(f:I\to\R\) 严格单调，令 \(J=f(I)\)。则 \(f:I\to J\) 为双射，并且逆函数
\[
 f^{-1}:J\to I
\]
严格单调，并且相对于子空间 \(J\) 连续。若再假设 \(f\) 连续，则 \(J\) 是区间。
\end{theorem}

\begin{proof}
严格单调性给出单射；按 \(J=f(I)\) 的定义得到满射。以下直接证明逆映射连续，不使用后文的介值定理。

设 \(f\) 严格递增，取 \(y_0=f(x_0)\in J\)。先设 \(x_0\) 是 \(I\) 的内点。给定 \(\varepsilon>0\)，选择 \(x_-,x_+\in I\) 使
\[
 x_0-\varepsilon<x_-<x_0<x_+<x_0+\varepsilon.
\]
严格递增给出
\[
 f(x_-)<y_0<f(x_+).
\]
令
\[
 \delta=\min\{y_0-f(x_-),\,f(x_+)-y_0\}>0.
\]
若 \(y\in J\) 且 \(\abs{y-y_0}<\delta\)，则
\[
 f(x_-)<y<f(x_+),
\]
从而 \(x_-<f^{-1}(y)<x_+\)，故
\[
 \abs{f^{-1}(y)-x_0}<\varepsilon.
\]
端点按相对邻域只选可用的一侧。递减情形把不等号方向反转即可。
\end{proof}

若 \(f\) 有跳跃，像集 \(J\) 可能有缺口，但上述相对连续性仍成立。只有在还要断言 \(J\) 为区间时，才需要 \(f\) 的连续性；这将在\chapterref{chap:continuous-functions-closed-interval}由介值定理推出。

\section*{习题}
\addcontentsline{toc}{section}{习题}

\subsection*{A组　定义与基本方法}
\begin{enumerate}[label=\textbf{\thechapter.\arabic*.}]
 \exerciseitem{17.1} 证明孤立定义域点处的任意函数都连续。
 \exerciseitem{17.2} 证明点连续与“极限等于点值”的等价性。
 \exerciseitem{17.3} 用量词否定构造数列，证明序列连续性判据的充分性。
 \exerciseitem{17.4} 证明连续函数保持绝对值、最大值与最小值运算。
 \exerciseitem{17.5} 证明若 \(f\) 在 \(a\) 连续且 \(f(a)\ne0\)，则 \(1/f\) 在 \(a\) 连续。
 \exerciseitem{17.6} 证明复合连续性，并分别用量词法和序列法完成。
 \exerciseitem{17.7} 证明 \(\sqrt x\) 在 \([0,\infty)\) 上连续，特别处理端点 \(0\)。
 \exerciseitem{17.8} 说明连续函数限制到任意子集后仍连续。
\end{enumerate}

\subsection*{B组　证明与反例}
\begin{enumerate}[label=\textbf{\thechapter.\arabic*.},resume]
 \exerciseitem{17.9} 对下列函数在指定点分类：\((x^2-1)/(x-1)\) 于 \(1\)，\(\lfloor x\rfloor\) 于整数，\(\sin(1/x)\) 于 \(0\)，\(1/x\) 于 \(0\)。
 \exerciseitem{17.10} 构造一个分段函数，使两段各自连续但交界点不连续；给出使其连续的参数条件。
 \exerciseitem{17.11} 证明递增函数的左极限等于左侧值集上确界。
 \exerciseitem{17.12} 写出递减函数的左右极限公式并证明。
 \exerciseitem{17.13} 证明单调函数的内点间断必为跳跃间断。
 \exerciseitem{17.14} 证明单调函数的跳跃区间两两不交。
 \exerciseitem{17.15} 完整证明单调函数间断点至多可数。
 \exerciseitem{17.16} 构造在可数无穷多个指定点发生跳跃的有界递增函数。
 \exerciseitem{17.17} 给出有可数稠密间断点集的单调函数。
 \exerciseitem{17.18} 严格单调函数若有一个跳跃，证明其像集不是区间。
\end{enumerate}

\subsection*{C组　综合问题}
\begin{enumerate}[label=\textbf{\thechapter.\arabic*.},resume]
 \exerciseitem{17.19} 证明严格单调连续函数的逆函数连续，单独说明区间端点。
 \exerciseitem{17.20} 证明若 \(f:I\to J\) 是严格递增双射且 \(I,J\) 都是区间，则 \(f\) 自动连续。
 \optionalexerciseitem{17.21} 研究“连续双射的逆函数连续”在非紧区间上何时成立，并与单调性联系。
 \optionalexerciseitem{17.22} 构造处处不连续但在每一点左右极限都存在的函数是否可能？给出证明。
 \projectexerciseitem{17.23} \ 【前置：\chapterref{chap:cauchy-construction}可数性；预计用时：75分钟】重建单调函数间断点至多可数的证明，并进一步证明跳跃大于 \(1/n\) 的间断点在有界值域内只有有限多个。
 \openexerciseitem{17.24} 比较可去、跳跃、第二类分类与函数振荡量
 \[
 \omega_f(a)=\inf_{\delta>0}\sup\{\abs{f(x)-f(y)}:x,y\in D,\ \abs{x-a},\abs{y-a}<\delta\}.
 \]
 \exerciseitem{17.25} 证明连续函数的零点集是闭集。
 \optionalexerciseitem{17.26} 设 \(f_n\) 都在 \(a\) 连续且 \(f_n\to f\) 一致。证明 \(f\) 在 \(a\) 连续；一致收敛的定义可先按\chapterref{chap:uniform-continuity}预读。
\end{enumerate}

